\documentclass[10pt]{article}
\usepackage[a4paper,margin=14mm,headheight=15pt]{geometry}
\usepackage{xcolor}
\usepackage{graphicx}
\usepackage{booktabs}
\usepackage{tabularx}
\usepackage{array}
\usepackage[hidelinks]{hyperref}

\definecolor{AgoraInk}{HTML}{17202A}
\definecolor{AgoraMuted}{HTML}{5D6D7E}
\definecolor{AgoraTeal}{HTML}{168A8A}
\definecolor{AgoraGold}{HTML}{D4A72C}
\definecolor{AgoraRed}{HTML}{C44E52}
\definecolor{AgoraPale}{HTML}{F1F5F5}
\definecolor{AgoraLine}{HTML}{CAD3D8}

\color{AgoraInk}
\setlength{\parindent}{0pt}
\setlength{\parskip}{4pt}
\setlength{\leftmargini}{4mm}
\makeatletter
\def\@listi{\leftmargin\leftmargini\parsep0pt\topsep2pt\itemsep1.5pt}
\makeatother
\renewcommand{\arraystretch}{1.28}
\pagestyle{myheadings}
\markright{Agora Human World Evaluation \hfill Version 2026-09-09\qquad}
\setlength{\fboxsep}{7pt}
\setlength{\fboxrule}{0.8pt}

\newcommand{\agoratitle}[2]{%
  {\fontsize{24}{27}\selectfont\bfseries #1\par}
  \vspace{2pt}
  {\large\color{AgoraTeal}\bfseries #2\par}
  \vspace{7pt}\textcolor{AgoraGold}{\rule{\linewidth}{2.2pt}}\vspace{4pt}
}

\newcommand{\worldtitle}[2]{%
  {\fontsize{17}{20}\selectfont\bfseries\textcolor{AgoraTeal}{#1}\quad #2\par}
  \vspace{3pt}\textcolor{AgoraGold}{\rule{\linewidth}{1.8pt}}\vspace{3pt}
}

\newcommand{\premisebox}[2]{%
  \noindent\fcolorbox{AgoraGold}{AgoraPale}{%
    \begin{minipage}{\dimexpr\linewidth-2\fboxsep-2\fboxrule\relax}
      \textbf{#1}\quad #2
    \end{minipage}}
}

\newcommand{\sectionlabel}[1]{%
  \par\vspace{3pt}{\large\bfseries\textcolor{AgoraTeal}{#1}}\par\vspace{-2pt}
}

\newcommand{\worldvisuals}[4]{%
  \begin{minipage}[t]{0.64\linewidth}
    \centering
    \includegraphics[width=\linewidth,height=0.34\textheight,keepaspectratio]{#1}\\[-1pt]
    {\small\color{AgoraMuted}#2}
  \end{minipage}\hfill
  \begin{minipage}[t]{0.32\linewidth}
    \centering
    \includegraphics[width=\linewidth,height=0.28\textheight,keepaspectratio]{#3}\\[-1pt]
    {\small\color{AgoraMuted}#4}
  \end{minipage}
}

\newcommand{\ratingheader}[2]{%
  \textbf{#1} & \textbf{1} & \textbf{2} & \textbf{3} & \textbf{4} & \textbf{5} & \textbf{#2}\\
}

\newcommand{\ratingrow}[2]{%
  \textbf{#1} #2 & $\square$ & $\square$ & $\square$ & $\square$ & $\square$ & \\\hline
}

\newcommand{\shortline}{\par\vspace{4pt}\textcolor{AgoraLine}{\rule{\linewidth}{0.45pt}}}

\newcommand{\responseblock}[2]{%
  \vspace{4pt}\textbf{#1}\par
  {\small\color{AgoraMuted}#2}\par
  \shortline\shortline\shortline
}

\newcommand{\evidencepath}{../figure_sources/input_snapshot/frontend/assets/generated}

\usepackage{amssymb}

\begin{document}
\agoratitle{Human Evaluation of Generated Worlds}{English questionnaire · self-contained artifact review}

\textbf{Purpose.} This questionnaire evaluates the quality of five worlds generated from five one-sentence premises. Everything needed for evaluation is printed in this document. Do not open or run Agora. Judge the generated world shown here, not agent behavior, runtime reliability, or the generation model.

\textbf{Respondent code:} \rule{42mm}{0.4pt}\hfill
\textbf{Assigned start code (A--E):} \rule{20mm}{0.4pt}

\sectionlabel{How to review}
Start with the world matching your assigned code, continue alphabetically, and wrap after E. Read the premise first, then inspect the equally sized evidence set: one map, one representative character sheet, rules, places, people, objects, and possible player actions. Complete that world's response page before continuing. Apply the same standard to every world and do not rank worlds against one another.

\sectionlabel{Rating scale}
\begin{tabularx}{\linewidth}{>{\bfseries}p{0.06\linewidth}X}
1 & Strongly disagree: the evidence contradicts the statement or provides almost no support.\\
2 & Disagree: important elements are generic, missing, or poorly connected.\\
3 & Mixed: the world is serviceable, but support is uneven or only partly developed.\\
4 & Agree: the evidence is specific, coherent, and supports the statement well.\\
5 & Strongly agree: the world is unusually specific, mutually reinforcing, and compelling.\\
\end{tabularx}

\sectionlabel{What the seven items mean}
The first six items assess premise realization, systemic coherence, society and character, spatial and material specificity, meaningful interaction possibilities, and visual coherence. Their equal-weight mean is the primary World Quality score. The seventh item is a separate overall judgment. ``Interaction possibilities'' means what this documented world appears to let a player attempt; it does not ask whether those actions succeeded in a runtime.

\vfill
\textcolor{AgoraMuted}{No model names, system traces, runtime captures, or performance scores are shown. Each world receives the same evidence fields and the same questions.}

% -----------------------------------------------------------------------------
\clearpage
\worldtitle{World A}{Clockwork Rain Conservatory}
\premisebox{One-sentence premise}{In a city-sized glass conservatory powered by mechanical rain, aristocratic houses and botanical guilds compete for control over patented seeds, weather control, and social status through trade, sabotage, and intricate legal battles.}

\vspace{7pt}
\worldvisuals
  {\evidencepath/world_asset_sets/creator_20260706_053251_1ad0b086_r001/world_map_source.png}
  {Generated world map}
  {\evidencepath/clockwork_rain_conservatory_main_01/creator_20260706_053251_1ad0b086_r001/raw_character_128.png}
  {Representative character sheet}

\begin{minipage}[t]{0.48\linewidth}
\sectionlabel{Rules and institutions}
\begin{itemize}
\item Reputation functions as currency; unsupported public accusations can ruin the accuser.
\item Inter-guild disputes go through the Adjudicator's Court rather than direct confrontation.
\item Knowledge of the Grand Mechanism is protected; unauthorized access is treason.
\end{itemize}
\sectionlabel{Activity loops}
Industrial sabotage; reputation games; patent litigation; seed and weather-device trade.
\sectionlabel{Places}
Grand Atrium Exchange; Valerius Prime Botanicals Showroom; Aqua-Regula Weather Systems Kiosk; Conservatory Central Seed Bank.
\end{minipage}\hfill
\begin{minipage}[t]{0.48\linewidth}
\sectionlabel{People and motives}
\begin{itemize}
\item Lysandra Valerius, disgraced heiress and botanical smuggler, seeks her family's Ever-Bloom patent.
\item Silas Vane, patent enforcer, hunts the missing seed patent and its smuggler.
\item Elara Thorne, weather-musician and debt collector, wants her Rain-Song patent restored.
\item Bramwell Loricatus, notary and rain-debt arbitrator, plans to seize the Valerius showroom.
\end{itemize}
\sectionlabel{Objects and possible actions}
Chrono-bloom seed; forged pollen permit; bottled monsoon cog; time-lattice vine cutting; rain-debt ledger. A player could sabotage a mechanism, uncover a secret, or file a formal grievance.
\end{minipage}

\clearpage
\worldtitle{World A}{Response sheet}
\small
\begin{tabularx}{\linewidth}{>{\raggedright\arraybackslash}X*{5}{>{\centering\arraybackslash}p{8mm}}p{24mm}}
\ratingheader{Statement}{Optional note}
\hline
\ratingrow{A1. Premise realization.}{The generated content develops the sentence into a distinctive world rather than merely restating it.}
\ratingrow{A2. Systemic coherence.}{Rules, institutions, resources, conflicts, and activity loops fit together causally.}
\ratingrow{A3. Society and character.}{The roles, motives, and tensions imply a believable community with more than interchangeable characters.}
\ratingrow{A4. Spatial/material specificity.}{Places and objects are concrete, recognizable, and tied to how the world works.}
\ratingrow{A5. Interaction possibilities.}{The world suggests multiple meaningful things a player could try beyond ordinary conversation.}
\ratingrow{A6. Visual coherence.}{The map and character asset are readable and belong to the same world described in the text.}
\ratingrow{A7. Overall judgment.}{Taken as a whole, this is a high-quality generated world that I would want to explore.}
\end{tabularx}
\normalsize
\responseblock{First action}{What is the first consequential action you would try in this world?}
\responseblock{Most distinctive element}{Which generated detail most strongly gives this world its own identity?}
\responseblock{One improvement}{Which single generated element most needs improvement?}

% -----------------------------------------------------------------------------
\clearpage
\worldtitle{World B}{Aurora Court of Migrating Cities}
\premisebox{One-sentence premise}{A convoy of nomadic, biomechanical cities traverses a frozen wasteland, their survival dependent on tense diplomacy, resource trading, and managing the complex logistics of their migration.}

\vspace{7pt}
\worldvisuals
  {\evidencepath/world_asset_sets/creator_20260724_203817_f86ae13e_r001/world_map_source.png}
  {Generated world map}
  {\evidencepath/aurora_court_of_migrating_cities_main_01/creator_20260724_203817_f86ae13e_r001/raw_character_128.png}
  {Representative character sheet}

\begin{minipage}[t]{0.48\linewidth}
\sectionlabel{Rules and institutions}
\begin{itemize}
\item Hospitality is a binding contract; breach triggers sanctions by the convoy.
\item A city's standing depends on the strength and wealth of its trade pacts.
\item The Aurora Court mediates formal disputes; unsanctioned conflict is piracy.
\end{itemize}
\sectionlabel{Activity loops}
Sanction cycles; resource scrambles; migration planning; pact formation.
\sectionlabel{Places}
Geothermal Conduit Hub; Lichen Vat Hydroponics; Cartographers' Spire; Scrap-Metal Exchange.
\end{minipage}\hfill
\begin{minipage}[t]{0.48\linewidth}
\sectionlabel{People and motives}
\begin{itemize}
\item Kaelen Voss, exchange arbitrator, needs rare-earth magnets for a levitation rail.
\item Elara Thorne, high cartographer, maps the Silent Zone and seeks a lost black box.
\item Thistle Vane, lichen-vat master, wants to prove biofuel can sustain the convoy.
\item Orlaith ``Oil-Slick'' Jax, chief siphonist, hides a fractured core while seeking a stabilizer.
\end{itemize}
\sectionlabel{Objects and possible actions}
Cryo-cilia actuator; gyro-stabilizer; sun-kelp ration; geothermic core shard; sub-glacial chart. A player could propose a sanction, forge a migration pact, or sponsor an expedition.
\end{minipage}

\clearpage
\worldtitle{World B}{Response sheet}
\small
\begin{tabularx}{\linewidth}{>{\raggedright\arraybackslash}X*{5}{>{\centering\arraybackslash}p{8mm}}p{24mm}}
\ratingheader{Statement}{Optional note}\hline
\ratingrow{B1. Premise realization.}{The generated content develops the sentence into a distinctive world rather than merely restating it.}
\ratingrow{B2. Systemic coherence.}{Rules, institutions, resources, conflicts, and activity loops fit together causally.}
\ratingrow{B3. Society and character.}{The roles, motives, and tensions imply a believable community with more than interchangeable characters.}
\ratingrow{B4. Spatial/material specificity.}{Places and objects are concrete, recognizable, and tied to how the world works.}
\ratingrow{B5. Interaction possibilities.}{The world suggests multiple meaningful things a player could try beyond ordinary conversation.}
\ratingrow{B6. Visual coherence.}{The map and character asset are readable and belong to the same world described in the text.}
\ratingrow{B7. Overall judgment.}{Taken as a whole, this is a high-quality generated world that I would want to explore.}
\end{tabularx}
\normalsize
\responseblock{First action}{What is the first consequential action you would try in this world?}
\responseblock{Most distinctive element}{Which generated detail most strongly gives this world its own identity?}
\responseblock{One improvement}{Which single generated element most needs improvement?}

% -----------------------------------------------------------------------------
\clearpage
\worldtitle{World C}{Mycelium Patent Bazaar}
\premisebox{One-sentence premise}{In a vast undercity built from engineered mycelium, rival guilds and families trade in living technology, from fungal bridges to encrypted spore data. The market thrives on innovation, intellectual property, and corporate espionage.}

\vspace{7pt}
\worldvisuals
  {\evidencepath/world_asset_sets/creator_20260724_211037_1c17ca96_r001/world_map_source.png}
  {Generated world map}
  {\evidencepath/mycelium_patent_bazaar_main_01/creator_20260724_211037_1c17ca96_r001/raw_character_128.png}
  {Representative character sheet}

\begin{minipage}[t]{0.48\linewidth}
\sectionlabel{Rules and institutions}
\begin{itemize}
\item Disputes use arbitration and sabotage; physical violence brings expulsion and asset seizure.
\item Reputation is currency, contracts are binding, and defaults become public.
\item Traders use identity masks; a true name creates legal and financial liability.
\end{itemize}
\sectionlabel{Activity loops}
Patent races; espionage and acquisition; biological grafting; injunctions.
\sectionlabel{Places}
Spore-Lock Depository; Graft-Wright's Stall; Patent Scribe's Office; Myco-Nostro Negotiation Lounge.
\end{minipage}\hfill
\begin{minipage}[t]{0.48\linewidth}
\sectionlabel{People and motives}
\begin{itemize}
\item Vespera Thorne, senior patent arbitrator, seeks the Ever-Bloom sequence.
\item Kaelen Voss, rogue bloom-broker, wants to register a forbidden neural-link patent.
\item Elara Vane, master graft-wright, aims to make the patent office obsolete.
\item Jaxen Krell, chitin logistics foreman, plans to intercept a primordial spore patent.
\end{itemize}
\sectionlabel{Objects and possible actions}
MycoCrypt spore key; lockpick fungus; Sun-Eater algae patent; mycelial patch; glimmer-cap vial. A player could verify a patent, initiate sabotage, or file an injunction.
\end{minipage}

\clearpage
\worldtitle{World C}{Response sheet}
\small
\begin{tabularx}{\linewidth}{>{\raggedright\arraybackslash}X*{5}{>{\centering\arraybackslash}p{8mm}}p{24mm}}
\ratingheader{Statement}{Optional note}\hline
\ratingrow{C1. Premise realization.}{The generated content develops the sentence into a distinctive world rather than merely restating it.}
\ratingrow{C2. Systemic coherence.}{Rules, institutions, resources, conflicts, and activity loops fit together causally.}
\ratingrow{C3. Society and character.}{The roles, motives, and tensions imply a believable community with more than interchangeable characters.}
\ratingrow{C4. Spatial/material specificity.}{Places and objects are concrete, recognizable, and tied to how the world works.}
\ratingrow{C5. Interaction possibilities.}{The world suggests multiple meaningful things a player could try beyond ordinary conversation.}
\ratingrow{C6. Visual coherence.}{The map and character asset are readable and belong to the same world described in the text.}
\ratingrow{C7. Overall judgment.}{Taken as a whole, this is a high-quality generated world that I would want to explore.}
\end{tabularx}
\normalsize
\responseblock{First action}{What is the first consequential action you would try in this world?}
\responseblock{Most distinctive element}{Which generated detail most strongly gives this world its own identity?}
\responseblock{One improvement}{Which single generated element most needs improvement?}

% -----------------------------------------------------------------------------
\clearpage
\worldtitle{World D}{Sunken Satellite Monastery}
\premisebox{One-sentence premise}{A monastic order salvages a crashed deep-space satellite, creating a unique underwater society where information, bandwidth, and salvaged tech are the primary currency.}

\vspace{7pt}
\worldvisuals
  {\evidencepath/world_asset_sets/creator_20260724_224121_43ca9360_r001/world_map_source.png}
  {Generated world map}
  {\evidencepath/sunken_satellite_monastery_main_01/creator_20260724_224121_43ca9360_r001/raw_character_128.png}
  {Representative character sheet}

\begin{minipage}[t]{0.48\linewidth}
\sectionlabel{Rules and institutions}
\begin{itemize}
\item Salvage claims must be registered with the Archivist before exchange.
\item Interrupting communion with the satellite data core is a grave offense.
\item Contracts encoded on data wafers bind access to information and machinery.
\end{itemize}
\sectionlabel{Activity loops}
Salvage cycles; information trade; signal interpretation; pressure-bound expeditions.
\sectionlabel{Places}
Data-Sanctum; Bandwidth Exchange; Scrap-Scribe's Workshop; Silent Canteen.
\end{minipage}\hfill
\begin{minipage}[t]{0.48\linewidth}
\sectionlabel{People and motives}
\begin{itemize}
\item Arch-Scribe Vaelen, void-frequency cantor, wants to decrypt the First Signal.
\item Kaelith-9, pressure scavenger and relic tuner, searches for the Origin-Key.
\item Sister Thalassa-Rho cultivates a biological filter against the signal.
\item Exarch Zephyrin-Mu, signal smuggler and bandwidth broker, plans to sell the intercepted frequency.
\end{itemize}
\sectionlabel{Objects and possible actions}
Gyroscopic prayer wheel; captured-starlight shard; encrypted Pater Noster; sonar canticle emitter; abyssal eye pearl. A player could decrypt a fragment, interface with the core, or negotiate a salvage contract.
\end{minipage}

\clearpage
\worldtitle{World D}{Response sheet}
\small
\begin{tabularx}{\linewidth}{>{\raggedright\arraybackslash}X*{5}{>{\centering\arraybackslash}p{8mm}}p{24mm}}
\ratingheader{Statement}{Optional note}\hline
\ratingrow{D1. Premise realization.}{The generated content develops the sentence into a distinctive world rather than merely restating it.}
\ratingrow{D2. Systemic coherence.}{Rules, institutions, resources, conflicts, and activity loops fit together causally.}
\ratingrow{D3. Society and character.}{The roles, motives, and tensions imply a believable community with more than interchangeable characters.}
\ratingrow{D4. Spatial/material specificity.}{Places and objects are concrete, recognizable, and tied to how the world works.}
\ratingrow{D5. Interaction possibilities.}{The world suggests multiple meaningful things a player could try beyond ordinary conversation.}
\ratingrow{D6. Visual coherence.}{The map and character asset are readable and belong to the same world described in the text.}
\ratingrow{D7. Overall judgment.}{Taken as a whole, this is a high-quality generated world that I would want to explore.}
\end{tabularx}
\normalsize
\responseblock{First action}{What is the first consequential action you would try in this world?}
\responseblock{Most distinctive element}{Which generated detail most strongly gives this world its own identity?}
\responseblock{One improvement}{Which single generated element most needs improvement?}

% -----------------------------------------------------------------------------
\clearpage
\worldtitle{World E}{Tidal Embassy of Lost Languages}
\premisebox{One-sentence premise}{A floating embassy is the sole hub for trading endangered languages, which are treated as political assets. Linguists, spies, and diplomats vie for control over translation rights and cultural memory.}

\vspace{7pt}
\worldvisuals
  {\evidencepath/world_asset_sets/creator_20260724_215709_f5f54ceb_r001/world_map_source.png}
  {Generated world map}
  {\evidencepath/tidal_embassy_of_lost_languages_main_01/creator_20260724_215709_f5f54ceb_r001/raw_character_128.png}
  {Representative character sheet}

\begin{minipage}[t]{0.48\linewidth}
\sectionlabel{Rules and institutions}
\begin{itemize}
\item Public emotion is a grave diplomatic breach inside the embassy.
\item Linguistic knowledge functions as currency, leverage, and political power.
\item Misusing a language is treated as a consequential political act.
\end{itemize}
\sectionlabel{Activity loops}
Whispered treason; ghosts in the protocol; translation auctions; provenance disputes.
\sectionlabel{Places}
Phoneme Exchange; Scriptorium of Whispers; Silent Auction Room; Cartographer's Nook.
\end{minipage}\hfill
\begin{minipage}[t]{0.48\linewidth}
\sectionlabel{People and motives}
\begin{itemize}
\item Elara Vane, phonetic smuggler and lexical restorationist, seeks the Siren's Syntax.
\item Valerius Thorne, lexical broker, wants a monopoly on the First Tongue.
\item Inspector Kaelen Voss investigates distribution of the Silence Plague.
\item Lyraen ``Ink-Lung'' Sunder forges memory maps to recover a family dialect.
\end{itemize}
\sectionlabel{Objects and possible actions}
Whispering conch; choral symbiote; Seal of the Silent Court; ghost-transcript slate; Ink of Echoes. A player could decode an echo, introduce a memetic phrase, or forge linguistic provenance.
\end{minipage}

\clearpage
\worldtitle{World E}{Response sheet}
\small
\begin{tabularx}{\linewidth}{>{\raggedright\arraybackslash}X*{5}{>{\centering\arraybackslash}p{8mm}}p{24mm}}
\ratingheader{Statement}{Optional note}\hline
\ratingrow{E1. Premise realization.}{The generated content develops the sentence into a distinctive world rather than merely restating it.}
\ratingrow{E2. Systemic coherence.}{Rules, institutions, resources, conflicts, and activity loops fit together causally.}
\ratingrow{E3. Society and character.}{The roles, motives, and tensions imply a believable community with more than interchangeable characters.}
\ratingrow{E4. Spatial/material specificity.}{Places and objects are concrete, recognizable, and tied to how the world works.}
\ratingrow{E5. Interaction possibilities.}{The world suggests multiple meaningful things a player could try beyond ordinary conversation.}
\ratingrow{E6. Visual coherence.}{The map and character asset are readable and belong to the same world described in the text.}
\ratingrow{E7. Overall judgment.}{Taken as a whole, this is a high-quality generated world that I would want to explore.}
\end{tabularx}
\normalsize
\responseblock{First action}{What is the first consequential action you would try in this world?}
\responseblock{Most distinctive element}{Which generated detail most strongly gives this world its own identity?}
\responseblock{One improvement}{Which single generated element most needs improvement?}

\end{document}
